\chapter{OOP Foundations and SOLID}
\chapterepigraph{Low-level design is the craft of turning a vague
requirement into a set of classes that are easy to change. Every design
pattern and every case study in this book is, at bottom, an application of
two things: the four pillars of object orientation, and the five SOLID
principles that keep those pillars from cracking as requirements evolve.
Learn these first and the rest of the book becomes recognition, not
memorisation.}

\section{Why This Chapter Comes First}
\label{sec:lld-intro}

Interviewers rarely ask ``recite the four pillars.'' They ask you to
\emph{design} something -- a parking lot, a notification service, a game --
and then watch whether your classes have clear responsibilities, whether
new features would force you to rewrite old code, and whether you can
justify each choice. Those judgements all trace back to the fundamentals
in this chapter.

\begin{ideabox}[The One Question Behind All Good Design]
Every design decision reduces to: \emph{``when the requirement changes,
how much code must I touch, and how likely am I to break something that
already works?''} Encapsulation, abstraction, inheritance and
polymorphism are the tools; the SOLID principles are the discipline that
aims those tools at \textbf{minimising the blast radius of change}. Keep
this question in mind and you can derive the right answer in an interview
instead of recalling one.
\end{ideabox}

\subsection*{How to Recognise Which Principle a Design Smell Violates}

\begin{patternbox}[Design Smells and the Principle They Break]
\begin{itemize}
  \item \textbf{A class that changes for many unrelated reasons}, or a
        ``manager''/``utils'' god class -- \textbf{Single Responsibility}.
  \item \textbf{Adding a new type forces editing a big
        \texttt{if/else} or \texttt{switch}} on a type field --
        \textbf{Open/Closed} (and usually a missing polymorphic
        hierarchy or a Factory/Strategy pattern).
  \item \textbf{A subclass that throws ``not supported'' or weakens a
        base-class guarantee} (the classic Square-extends-Rectangle) --
        \textbf{Liskov Substitution}.
  \item \textbf{A fat interface whose implementers are forced to stub out
        methods they do not need} -- \textbf{Interface Segregation}.
  \item \textbf{High-level policy code that \texttt{new}s concrete
        low-level classes directly}, making it impossible to swap
        implementations or unit-test -- \textbf{Dependency Inversion}
        (inject an abstraction instead).
\end{itemize}
\end{patternbox}

The three sections of this chapter cover the four pillars, the ways
classes relate to one another (association, aggregation, composition, and
inheritance, with their UML notation), and the SOLID principles with a
concrete before/after refactor for each.

\newpage

\section{The Four Pillars of Object Orientation}
\label{sec:lld-four-pillars}

\problemstatement{Explain encapsulation, abstraction, inheritance, and
polymorphism -- what each means, why it exists, and how it looks in C++ --
since every later pattern is built from them.}

\begin{patternbox}[What Each Pillar Is Really For]
\begin{itemize}
  \item \textbf{Encapsulation} -- bundle data with the methods that guard
        it, and hide the data behind those methods, so invariants cannot
        be broken from outside.
  \item \textbf{Abstraction} -- expose \emph{what} an object does, hide
        \emph{how}, so callers depend on a stable interface, not
        internals.
  \item \textbf{Inheritance} -- reuse and specialise a base type,
        modelling an ``is-a'' relationship.
  \item \textbf{Polymorphism} -- let one interface stand for many
        concrete behaviours chosen at run time, so calling code needn't
        know the exact subtype.
\end{itemize}
\end{patternbox}

\subsection*{Encapsulation and abstraction: two sides of the same wall}

Encapsulation is about \emph{protection}: making fields \texttt{private}
and offering controlled access so no caller can leave an object in an
invalid state (a bank balance that only \texttt{deposit}/\texttt{withdraw}
can change, never a raw assignment). Abstraction is about
\emph{simplification}: the caller sees \texttt{account.withdraw(100)} and
neither knows nor cares how overdraft rules are enforced. Encapsulation
hides data; abstraction hides implementation. In practice the same
\texttt{private} fields plus a small public interface deliver both.

\begin{ideabox}[Program to an Interface, Not an Implementation]
The single most valuable habit the four pillars enable: callers should
depend on an \emph{abstract type} (a base class or interface), so the
concrete class behind it can be swapped without touching the caller. This
one idea is the seed of polymorphism, the Strategy pattern, dependency
injection, and testability.
\end{ideabox}

\subsection*{Inheritance and polymorphism: one call, many behaviours}

Inheritance lets a \texttt{Dog} reuse and extend an \texttt{Animal}.
Polymorphism lets a function that accepts an \texttt{Animal*} call
\texttt{speak()} and get barking or meowing depending on the real object,
resolved at run time through a \texttt{virtual} function. This run-time
dispatch is what allows open-ended designs: new animal types slot in
without changing the function that iterates over animals.

\subsection*{C++ illustration}

\begin{lstlisting}
#include <bits/stdc++.h>
using namespace std;

// Encapsulation + abstraction: data hidden, invariant enforced.
class Account {
    double balance = 0;                 // private: cannot be corrupted
public:
    void deposit(double amt) { if (amt > 0) balance += amt; }
    bool withdraw(double amt) {
        if (amt > 0 && amt <= balance) { balance -= amt; return true; }
        return false;                   // overdraft rule hidden inside
    }
    double getBalance() const { return balance; }
};

// Inheritance + polymorphism: one interface, many behaviours.
class Animal {
public:
    virtual string speak() const = 0;   // pure virtual: abstraction
    virtual ~Animal() = default;
};
class Dog : public Animal {
public:
    string speak() const override { return "Woof"; }
};
class Cat : public Animal {
public:
    string speak() const override { return "Meow"; }
};

int main() {
    Account a;
    a.deposit(100);
    cout << "Withdrew 40? " << a.withdraw(40) << ", balance "
         << a.getBalance() << "\n";

    vector<unique_ptr<Animal>> zoo;
    zoo.push_back(make_unique<Dog>());
    zoo.push_back(make_unique<Cat>());
    for (auto& animal : zoo)            // caller knows only "Animal"
        cout << animal->speak() << " ";
    cout << "\n";
    return 0;
}
\end{lstlisting}

\begin{takeawaybox}
Encapsulation guards invariants, abstraction hides implementation,
inheritance reuses, and polymorphism defers the choice of behaviour to run
time. The recurring pay-off is the same: \emph{callers depend on a stable
abstraction, so concrete details can change without a ripple}. Every
pattern in the next chapter is a specific, named way of exploiting that.
\end{takeawaybox}

\section{Class Relationships and UML Notation}
\label{sec:lld-relationships}

\problemstatement{Distinguish the ways two classes can relate --
association, aggregation, composition, dependency, and inheritance -- and
know the UML class-diagram notation, since interviewers expect you to
sketch a design and name these relationships correctly.}

\begin{patternbox}[Which Relationship Is It?]
\begin{itemize}
  \item \textbf{``is-a''} -- \textbf{inheritance} (a Car is-a Vehicle).
        UML: hollow triangle arrow to the parent.
  \item \textbf{``has-a,'' parts can live independently} --
        \textbf{aggregation} (a Team has Players; players outlive the
        team). UML: hollow diamond.
  \item \textbf{``has-a,'' parts die with the whole} --
        \textbf{composition} (a House has Rooms; destroy the house and
        the rooms go with it). UML: filled diamond.
  \item \textbf{``uses-a'' transiently} -- \textbf{dependency} (an
        \texttt{Order} uses a \texttt{PaymentGateway} passed to one
        method). UML: dashed arrow.
  \item \textbf{A general structural link} -- \textbf{association}
        (a plain solid line); aggregation and composition are its
        stronger, ownership-bearing forms.
\end{itemize}
\end{patternbox}

\subsection*{The key axis: ownership and lifetime}

The distinction interviewers probe is \textbf{aggregation vs.\
composition}, and it is entirely about \emph{lifetime}. In composition the
whole \emph{owns} its parts and controls their creation and destruction --
model it by storing parts \emph{by value} or via an owning
\texttt{unique\_ptr}. In aggregation the whole merely \emph{references}
parts that exist independently -- model it by storing a non-owning pointer
or reference handed in from outside.

\begin{ideabox}[Favour Composition Over Inheritance]
A foundational LLD maxim. Inheritance is rigid: it fixes behaviour at
compile time and leaks the base class's details into every subclass. Composing
an object out of smaller collaborators (holding a \texttt{Strategy}, an
\texttt{Engine}, a \texttt{Logger}) lets you change behaviour by swapping a
part at run time. When tempted to subclass, first ask whether ``has-a''
models the problem better than ``is-a.''
\end{ideabox}

\subsection*{C++ illustration of each relationship}

\begin{lstlisting}
#include <bits/stdc++.h>
using namespace std;

class Engine {                 // an independent part
public:
    void start() { cout << "Engine started\n"; }
};

class Logger {                 // dependency: used transiently
public:
    void log(const string& m) { cout << "[log] " << m << "\n"; }
};

// Composition: Car owns its Engine (created and destroyed with the Car).
class Car {
    Engine engine;                     // by value -> composition
public:
    void drive(Logger& logger) {       // dependency: Logger passed in
        engine.start();
        logger.log("car driving");
    }
};

// Aggregation: Team references Players it does not own.
class Player {
public:
    string name;
    Player(string n) : name(move(n)) {}
};
class Team {
    vector<Player*> players;           // non-owning -> aggregation
public:
    void add(Player* p) { players.push_back(p); }
    void roster() {
        for (auto* p : players) cout << p->name << " ";
        cout << "\n";
    }
};

int main() {
    Car car; Logger logger;
    car.drive(logger);

    Player a("Alice"), b("Bob");       // players outlive any team
    Team t; t.add(&a); t.add(&b);
    t.roster();
    return 0;
}
\end{lstlisting}

\begin{takeawaybox}
Name relationships precisely: inheritance is ``is-a''; composition is
owned ``has-a'' (filled diamond, shared lifetime); aggregation is
referenced ``has-a'' (hollow diamond, independent lifetime); dependency is
a transient ``uses-a.'' In code the tell is ownership -- by-value/owning
pointer for composition, non-owning pointer for aggregation. When in
doubt, prefer composition over inheritance.
\end{takeawaybox}

\section{The SOLID Principles}
\label{sec:lld-solid}

\problemstatement{State each of the five SOLID principles, the design
smell it cures, and a concrete before/after in C++. These principles are
the vocabulary interviewers use to critique a design, so each is presented
as a problem it solves.}

\begin{patternbox}[SOLID in One Line Each]
\textbf{S}ingle Responsibility -- one reason to change per class.
\textbf{O}pen/Closed -- open to extension, closed to modification.
\textbf{L}iskov Substitution -- a subtype must be usable anywhere its base
is. \textbf{I}nterface Segregation -- many small interfaces beat one fat
one. \textbf{D}ependency Inversion -- depend on abstractions, not
concretions.
\end{patternbox}

\subsection*{S -- Single Responsibility Principle}

A class should have exactly one reason to change. An \texttt{Invoice} that
computes totals \emph{and} formats itself as PDF \emph{and} emails itself
has three reasons to change (tax rules, layout, mail server) and three
teams stepping on each other. Split it: \texttt{Invoice} (data + totals),
\texttt{InvoicePrinter} (formatting), \texttt{InvoiceMailer} (sending).

\subsection*{O -- Open/Closed Principle}

Software entities should be open for extension but closed for
modification. The smell is a growing \texttt{switch} on a type tag:

\begin{lstlisting}
// BAD: every new shape edits this function (closed to extension).
double area(const Shape& s) {
    if (s.type == CIRCLE)    return 3.14159 * s.r * s.r;
    if (s.type == RECTANGLE) return s.w * s.h;
    // ... adding Triangle means editing here again
}
\end{lstlisting}

Replace the tag with polymorphism so new shapes are \emph{added}, not
\emph{edited} in:

\begin{lstlisting}
struct Shape { virtual double area() const = 0; virtual ~Shape() = default; };
struct Circle : Shape {
    double r; Circle(double r_) : r(r_) {}
    double area() const override { return 3.14159 * r * r; }
};
struct Rectangle : Shape {
    double w, h; Rectangle(double w_, double h_) : w(w_), h(h_) {}
    double area() const override { return w * h; }
};
// A new Triangle class needs zero changes to existing code.
\end{lstlisting}

\subsection*{L -- Liskov Substitution Principle}

A subclass must honour every promise of its base, so code using the base
never breaks when handed the subclass. The classic violation: a
\texttt{Square} that inherits \texttt{Rectangle} and overrides
\texttt{setWidth} to also change the height silently breaks any code that
sets width and height independently and expects them to stay independent.
The fix is usually to \emph{not} force the inheritance -- model
\texttt{Square} and \texttt{Rectangle} as siblings under a common
\texttt{Shape}, or use composition.

\begin{ideabox}[LSP Is About Behaviour, Not Just Signatures]
Compiling is not enough. A subtype substitutes correctly only if it keeps
the base's \emph{behavioural contract}: it must not strengthen
preconditions, weaken postconditions, or throw ``unsupported'' for methods
the base guarantees. If an override has to disable a base method, the
inheritance is wrong.
\end{ideabox}

\subsection*{I -- Interface Segregation Principle}

No client should be forced to depend on methods it does not use. A fat
\texttt{Machine} interface with \texttt{print()}, \texttt{scan()},
\texttt{fax()} forces a simple printer to stub out \texttt{scan} and
\texttt{fax}. Split into \texttt{Printer}, \texttt{Scanner},
\texttt{Fax}; a multifunction device implements all three, a basic printer
implements only one.

\subsection*{D -- Dependency Inversion Principle}

High-level modules should depend on abstractions, not on concrete
low-level modules -- and the abstraction is injected, not constructed
inside:

\begin{lstlisting}
#include <bits/stdc++.h>
using namespace std;

struct MessageSender {                 // abstraction
    virtual void send(const string& msg) = 0;
    virtual ~MessageSender() = default;
};
struct EmailSender : MessageSender {
    void send(const string& msg) override { cout << "Email: " << msg << "\n"; }
};
struct SmsSender : MessageSender {
    void send(const string& msg) override { cout << "SMS: " << msg << "\n"; }
};

class NotificationService {            // high-level policy
    MessageSender& sender;             // depends on the abstraction
public:
    NotificationService(MessageSender& s) : sender(s) {}   // injected
    void notify(const string& msg) { sender.send(msg); }
};

int main() {
    EmailSender email; SmsSender sms;
    NotificationService viaEmail(email), viaSms(sms);
    viaEmail.notify("Order shipped");
    viaSms.notify("OTP 4321");         // swap implementation, no edits
    return 0;
}
\end{lstlisting}

\begin{takeawaybox}
SOLID is one theme in five costumes: \emph{isolate what changes}. SRP
isolates responsibilities, OCP isolates new behaviour behind polymorphism,
LSP keeps subtypes safe to substitute, ISP isolates clients from
irrelevant methods, and DIP isolates high-level policy from low-level
detail via an injected abstraction. Every pattern in the next chapter is a
reusable structure that upholds one or more of these.
\end{takeawaybox}

